\documentclass[aspectratio=169,10pt,spanish]{beamer}

\input{../../../_preambulo_beamer.tex}
\input{../../../_comandos_beamer.tex}

\selectlanguage{spanish}

\title{MA1001B - Modelación Estadística para la Toma de Decisiones}
\subtitle{Lección 24: Teorema del límite central}
\author{}
\institute{Tecnol\'ogico de Monterrey}
\date{}

\begin{document}

\begin{frame}[plain]
  \vspace{1.2cm}
  {\Large\bfseries MA1001B - Modelación Estadística para la Toma de Decisiones\par}
  \vspace{0.7cm}
  {\large Lección 24: Teorema del límite central\par}
  \vspace{0.7cm}
  {\color{TecRojo}\rule{\textwidth}{1.2pt}}
  \vspace{0.8cm}

  Facultad MA1001B\\[0.3cm]
  Periodo\\[0.3cm]
  Tecnol\'ogico de Monterrey
\end{frame}

\begin{frame}{La pregunta de la población asimétrica}
  \begin{alertblock}{Pregunta guía}
    Si los tiempos de retraso son exponenciales y no acampanados, ¿cuándo
    sigue siendo aproximadamente normal la distribución muestral de $\bar x$?
  \end{alertblock}

  \vspace{0.6em}
  Al final de esta lección, usted puede:
  \begin{itemize}
    \item reconocer una población sesgada a la derecha que no es normal;
    \item comparar histogramas de $\bar x$ para $n=4$ y $n=40$;
    \item enunciar el TLC para la media muestral;
    \item explicar por qué una cola a dos EE falla con $n$ pequeño.
  \end{itemize}

  \vfill
  \scriptsize Todos los tiempos de retraso usados hoy son sintéticos. No se usan datos reales de empresa.
\end{frame}

\begin{frame}{Un reloj exponencial de retrasos}
  \small
  $X\sim\mathrm{Exponencial}(\text{media}=10)$:
  \[
    E[X]=\sigma=10,\qquad \text{asimetría}=2,\qquad \text{mediana}=6.931472.
  \]
  La cola derecha es larga: $P(X>20)=0.135335$.
  Para muestras iid,
  \[
    E[\bar x]=10,\qquad \mathrm{EE}(\bar x)=10/\sqrt{n}.
  \]

  \vspace{0.3em}
  \begin{exampleblock}{Objetivo del TLC}
    El TLC es un enunciado sobre $\bar x$, no sobre que los retrasos
    originales se vuelvan normales.
  \end{exampleblock}
\end{frame}

\begin{frame}[fragile]{Paso 1: Población asimétrica}
  \begin{columns}[T]
    \begin{column}{0.42\textwidth}
      \small
      \begin{lstlisting}
model = expon(scale=10)
mean = model.mean()
skew = model.stats(
  moments="s"
)
      \end{lstlisting}

      \begin{block}{Salida verificada}
        Media $=$ sd $=10$\\
        Mediana $=6.931472$\\
        Asimetría $=2.000000$
      \end{block}
    \end{column}
    \begin{column}{0.58\textwidth}
      \centering
      \includegraphics[width=\textwidth]{../figuras/clt_01_skewed_population.png}
      \vspace{0.2em}
      \scriptsize Densidad exponencial de retrasos del Paso 1
    \end{column}
  \end{columns}
\end{frame}

\begin{frame}{Un $n$ mayor acerca $\bar x$ a una campana}
  \small
  5{,}000 muestras con semilla 42:
  \[
    n=4:\ \widehat{\mathrm{EE}}=5.076401,\ \text{asimetr\'ia}=1.049897.
  \]
  \[
    n=40:\ \widehat{\mathrm{EE}}=1.574371,\ \text{asimetr\'ia}=0.332572.
  \]
  La asimetría teórica de una media exponencial es $2/\sqrt{n}$. Ese valor
  es $1$ en $n=4$ y $0.316228$ en $n=40$.
\end{frame}

\begin{frame}[fragile]{Paso 2: $n=4$ versus $n=40$}
  \begin{columns}[T]
    \begin{column}{0.40\textwidth}
      \small
      \begin{lstlisting}
draws = model.rvs(
  size=(5000, n)
)
xbars = draws.mean(axis=1)
      \end{lstlisting}

      \begin{block}{Salida verificada}
        $n=4$ asimetría $=1.049897$\\
        $n=40$ asimetría $=0.332572$
      \end{block}
    \end{column}
    \begin{column}{0.60\textwidth}
      \centering
      \includegraphics[width=\textwidth]{../figuras/clt_02_means_n4_n40.png}
      \vspace{0.2em}
      \scriptsize Histogramas muestrales del Paso 2
    \end{column}
  \end{columns}
\end{frame}

\begin{frame}{Paso 3: Un $n$ pequeño hace fallar la cola}
  \begin{columns}[T]
    \begin{column}{0.42\textwidth}
      \small
      Ley exacta: $\bar x\sim\mathrm{Gamma}$.
      Evento: $\bar x>\mu+2\,\mathrm{EE}$.

      \vspace{0.4em}
      $n=4$: exacta $0.042380$ versus normal $0.022750$.\\
      $n=40$: exacta $0.030560$ versus normal $0.022750$.

      \vspace{0.4em}
      \textbf{Límite:} $n$ pequeño más asimetría hace que el TLC sea una mala
      aproximación. La Lección 25 estudia $\hat p$.
    \end{column}
    \begin{column}{0.58\textwidth}
      \centering
      \includegraphics[width=\textwidth]{../figuras/clt_03_small_n_skew_limit.png}
      \vspace{0.2em}
      \scriptsize Colas a dos EE exactas versus normales del Paso 3
    \end{column}
  \end{columns}
\end{frame}

\begin{frame}{Verificación de conceptos}
  \begin{enumerate}
    \item ¿Por qué la población de retrasos no es un modelo normal?
    \item ¿Qué ocurre con la asimetría de $\bar x$ cuando $n$ crece?
    \item En $n=4$, compare la cola exacta a dos EE con $0.022750$.
    \item ¿Hace $n=40$ que la cola normal sea exacta?
  \end{enumerate}

  \vspace{0.6em}
  \begin{block}{Conexión con la Lección 25}
    La sesión puede continuar directamente con la distribución muestral de
    una proporción muestral $\hat p$.
  \end{block}

  \vfill
  \scriptsize Fuente: OpenStax, \textit{Introductory Business Statistics 2e},
  Capítulo 7, Secciones 7.1--7.2. CC BY-NC-SA.
\end{frame}

\appendix



\end{document}
